===============================================================================
\item 塩化ナトリウム(\ce{NaCl}，58.44\ g/mol)\cite{7}\\
    \hspace{0pt}・常温常圧下では，白色の結晶〜結晶性粉末である。\\
    \hspace{0pt}・水に溶けやすく，エタノールには溶けにくい。\\
    
\bibitem{7}  厚生労働省：金塩化ナトリウム，職場のあんぜんサイト(オンライン)，\\
           入手先\url{<https://anzeninfo.mhlw.go.jp/anzen/gmsds/15189-51-2.html>}
           (参照2025-4-18).
===============================================================================
\item 無水硫酸マグネシウム(\ce{MgSO4}，120.37\ g/mol)\cite{8}\\
    \hspace{0pt}・常温常圧下では，白色の粉末であり，吸湿性を持つ。\\
    \hspace{0pt}・水に溶けやすく，エタノールには溶けにくい。\\
    
\bibitem{8}  富士フイルム 和光純薬株式会社：硫酸マグネシウム，富士フイルム 和光純薬株式会社(オンライン)，\\
           入手先\url{<https://labchem-wako.fujifilm.com/jp/product/detail/W01W0113-1233.html>}
           (参照2025-4-22).
===============================================================================
\item 水酸化カリウム(\ce{KOH}，56.11\ g/mol)\cite{9}\\
    \hspace{0pt}・常温常圧下では，白色の個体であり，潮解性を持つ。\\
    \hspace{0pt}・水に極めて溶けやすく，エタノールに溶けやすい。\\
    \hspace{0pt}・溶解時に激しい発熱をする。\\
    \hspace{0pt}・1価の強塩基性の物質である。\\

\bibitem{9}  厚生労働省：水酸化カリウム，職場のあんぜんサイト(オンライン)，\\
           入手先\url{<https://anzeninfo.mhlw.go.jp/anzen/gmsds/1310-58-3.html>}
           (参照2025-4-18).
===============================================================================
\item 濃硫酸(\ce{H2SO4}，\ g/mol)\cite{7}\\
                    \hspace{0pt}・常温常圧下では，無色透明の液体であり，強い脱水性がある。\\
                    \hspace{0pt}・水，アルコール及びエーテルと激しく発熱しながら任意の割合で混ざる。\\
                    \hspace{0pt}・強力な酸化剤であり，可燃性物質や還元性物質と反応する。\\
                    \hspace{0pt}・弱酸性の物質だが，硫酸の濃度が低くなると強酸性の物質である。\\
\bibitem{7}  厚生労働省:硫酸，職場のあんぜんサイト(オンライン)，\\
                         入手先\url{<https://anzeninfo.mhlw.go.jp/anzen/gmsds/0626.html>}
                         (参照2025-5-8).
===============================================================================                         
\item 塩化第二鉄(\ce{FeCl3}，162.22\ g/mol)\cite{12}\\
                    \hspace{0pt}・常温常圧下では，暗褐色の透明〜微濁の液体である。\\
                    \hspace{0pt}・水とアルコールに極めて溶けやすい。\\
                    \hspace{0pt}・200${}^\circ$C以上に加熱すると分解し，有毒で腐食性の気体が生じる。\\
                    \hspace{0pt}・水溶液は中程度の強さの酸性物質である。\\
\bibitem{12}  厚生労働省:塩化鉄\RRnumber{3}，職場のあんぜんサイト(オンライン)，\\
                         入手先\url{<https://anzeninfo.mhlw.go.jp/anzen/gmsds/7705-08-0.html>}
                         (参照2025-5-11).
===============================================================================
\item 塩酸(\ce{HCl}，36.46\ g/mol)\cite{4}\\
                    \hspace{0pt}・常温常圧下では，無色の液体であり，強い刺激臭がある。\\
                    \hspace{0pt}・強酸性の物質であり，多くの金属に対して腐食性を示す。\\
                    \hspace{0pt}・塩基性物質や酸化剤と激しく反応する。\\
\bibitem{4}  厚生労働省:塩化水素，職場のあんぜんサイト(オンライン)，\\
                         入手先\url{<https://anzeninfo.mhlw.go.jp/anzen/gmsds/7647-01-0.html>}
                         (参照2025-5-18).
===============================================================================
\item アンモニア水(\ce{NH3}，17.03\ g/mol)\cite{7}\\
                    \hspace{0pt}・常温常圧下では，無色透明の液体であり，特異臭がある。\\
                    \hspace{0pt}・水とエタノールに任意の割合で混和する。\\
                    \hspace{0pt}・強酸化剤やハロゲンと激しく反応する。\\
                    \hspace{0pt}・強塩基性の物質である。\\
\bibitem{7}  厚生労働省:アンモニア，職場のあんぜんサイト(オンライン)，\\
                         入手先\url{<https://anzeninfo.mhlw.go.jp/anzen/gmsds/7664-41-7.html>}
                         (参照2025-5-25).


\item 塩化アンモニウム(\ce{NH4Cl}，53.49\ g/mol)\cite{8}\\
                    \hspace{0pt}・常温常圧下では，無色〜白色の結晶である。\\
                    \hspace{0pt}・水に溶けやすく，エタノールに溶けにくい。\\
                    \hspace{0pt}・水溶液は弱酸性の物質である。\\
\bibitem{8}  厚生労働省:塩化アンモニウム，職場のあんぜんサイト(オンライン)，\\
                         入手先\url{<https://anzeninfo.mhlw.go.jp/anzen/gmsds/12125-02-9.html>}
                         (参照2025-5-25).  

























